% GENERATED FILE. Edit canonical lessons, then run npm run generate:books.
% canonical-source-hash: fnv1a64:f520b4a9711bb04b
% canonical-lessons: AR-C495-inziaj, AR-C495-munzaij, AR-C495-musta, AR-C495-tahammala, AR-C495-dajira, AR-R495-review-488-495

\chapter{Annoyed and fed up}
\label{ch:ar-annoyed-and-fed-up}
\hlchaptermodality{495}

\begin{chapteropening}
\textbf{By the end of this chapter:} \emph{I can talk about annoyance, say that someone is annoyed or displeased, and say that someone put up with something or was fed up.}

Complete the last lesson of chapter 495: I can talk about annoyance, say that someone is annoyed or displeased, and say that someone put up with something or was fed up.
\end{chapteropening}

\section[\texorpdfstring{\ar{انزعاج} (inziʿāj)}{انزعاج (inziʿāj)}]{\ar{انزعاج} (inziʿāj) — annoyance}

\label{lesson:AR-C495-inziaj}



\begin{quote}
Before the new one: say the Arabic for frustrated, then the Arabic for to get bored.
\end{quote}

\subsection*{You'll want to know: \ar{انزعاج}}
\textbf{\ar{انزعاج}} — \emph{inziʿāj} — \textquotedblleft{}annoyance\textquotedblright{}.

\begin{quote}
\ar{تكلم} \ar{الجار} \ar{عن} \ar{الموسيقى} \ar{العالية} \ar{بانزعاج}.
\end{quote}

\begin{quote}
\emph{takallama al-jār ʿan al-mūsīqā al-ʿāliya bi-inziʿāj.}
\end{quote}

\begin{quote}
The neighbour spoke about the loud music with annoyance.
\end{quote}

The adjective \ar{مزعج}, annoying, is from the same root.

\begin{cousinweb}
\textbf{\ar{انزعاج}} (\emph{inziʿāj}) is built on the native Arabic root \emph{z-ʿ-j}, \textquotedblleft{}to disquiet, to unsettle\textquotedblright{} (Lane). Inziʿāj is the verbal noun of Form VII: being unsettled.
\end{cousinweb}

\begin{grammarlens}[title={What you've built}]
One more word for describing what you lived and felt.
\end{grammarlens}

\subsection*{Your turn}
\begin{itemize}
\raggedright
  \item \emph{Say it:} \emph{inziʿāj}
  \item \emph{Say it:} \emph{inziʿāj}, once more
  \item \emph{Say it:} \emph{takallama al-jār ʿan al-mūsīqā al-ʿāliya bi-inziʿāj.} — The neighbour spoke about the loud music with annoyance.
  \item \emph{Recall:} say the Arabic for frustrated, then the Arabic for to get bored, then say \emph{inziʿāj} again
\end{itemize}

\begin{tcolorbox}[breakable,colback=teal!4,colframe=teal!35!black,title={Before you move on}]
What does \ar{انزعاج} mean? (\textquotedblleft{}Annoyance\textquotedblright{}.) Say it once more.
\end{tcolorbox}
\section[\texorpdfstring{\ar{منزعج} (munzaʿij)}{منزعج (munzaʿij)}]{\ar{منزعج} (munzaʿij) — annoyed}

\label{lesson:AR-C495-munzaij}



\begin{quote}
Before the new one: say the Arabic for to get bored, then the Arabic for annoyance.
\end{quote}

\subsection*{You'll want to know: \ar{منزعج}}
\textbf{\ar{منزعج}} — \emph{munzaʿij} — \textquotedblleft{}annoyed\textquotedblright{}.

\begin{quote}
\ar{الأب} \ar{منزعج} \ar{من} \ar{صوت} \ar{التلفاز}.
\end{quote}

\begin{quote}
\emph{al-ab munzaʿij min ṣawt at-tilfāz.}
\end{quote}

\begin{quote}
The father is annoyed by the noise of the television.
\end{quote}

It takes \ar{من} for the cause.

\begin{cousinweb}
\textbf{\ar{منزعج}} (\emph{munzaʿij}) is built on the native Arabic root \emph{z-ʿ-j}, \textquotedblleft{}to disquiet, to unsettle\textquotedblright{} (Lane). Munzaʿij is the participle of Form VII.
\end{cousinweb}

\begin{grammarlens}[title={What you've built}]
One more word for describing what you lived and felt.
\end{grammarlens}

\subsection*{Your turn}
\begin{itemize}
\raggedright
  \item \emph{Say it:} \emph{munzaʿij}
  \item \emph{Say it:} \emph{munzaʿij}, once more
  \item \emph{Say it:} \emph{al-ab munzaʿij min ṣawt at-tilfāz.} — The father is annoyed by the noise of the television.
  \item \emph{Recall:} say the Arabic for to get bored, then the Arabic for annoyance, then say \emph{munzaʿij} again
\end{itemize}

\begin{tcolorbox}[breakable,colback=teal!4,colframe=teal!35!black,title={Before you move on}]
What does \ar{منزعج} mean? (\textquotedblleft{}Annoyed\textquotedblright{}.) Say it once more.
\end{tcolorbox}
\section[\texorpdfstring{\ar{مستاء} (mustāʾ)}{مستاء (mustāʾ)}]{\ar{مستاء} (mustāʾ) — displeased, upset}

\label{lesson:AR-C495-musta}



\begin{quote}
Before the new one: say the Arabic for annoyance, then the Arabic for annoyed.
\end{quote}

\subsection*{You'll want to know: \ar{مستاء}}
\textbf{\ar{مستاء}} — \emph{mustāʾ} — \textquotedblleft{}displeased, upset\textquotedblright{}.

\begin{quote}
\ar{الزبون} \ar{مستاء} \ar{من} \ar{الخدمة}.
\end{quote}

\begin{quote}
\emph{az-zabūn mustāʾ min al-khidma.}
\end{quote}

\begin{quote}
The customer is unhappy with the service.
\end{quote}

\begin{cousinweb}
\textbf{\ar{مستاء}} (\emph{mustāʾ}) is built on the native Arabic root \emph{s-w-ʾ}, \textquotedblleft{}to be evil, to be bad\textquotedblright{} (Lane). Mustāʾ is the participle of Form VIII, istāʾa: one who takes a thing badly.
\end{cousinweb}

\begin{grammarlens}[title={What you've built}]
One more word for describing what you lived and felt.
\end{grammarlens}

\subsection*{Your turn}
\begin{itemize}
\raggedright
  \item \emph{Say it:} \emph{mustāʾ}
  \item \emph{Say it:} \emph{mustāʾ}, once more
  \item \emph{Say it:} \emph{az-zabūn mustāʾ min al-khidma.} — The customer is unhappy with the service.
  \item \emph{Recall:} say the Arabic for annoyance, then the Arabic for annoyed, then say \emph{mustāʾ} again
\end{itemize}

\begin{tcolorbox}[breakable,colback=teal!4,colframe=teal!35!black,title={Before you move on}]
What does \ar{مستاء} mean? (\textquotedblleft{}Displeased, upset\textquotedblright{}.) Say it once more.
\end{tcolorbox}
\section[\texorpdfstring{\ar{تحمل} (taḥammala)}{تحمل (taḥammala)}]{\ar{تحمل} (taḥammala) — to put up with}

\label{lesson:AR-C495-tahammala}



\begin{quote}
Before the new one: say the Arabic for annoyed, then the Arabic for displeased.
\end{quote}

\subsection*{You'll want to know: \ar{تحمل}}
\textbf{\ar{تحمل}} — \emph{taḥammala} — \textquotedblleft{}to put up with\textquotedblright{}.

\begin{quote}
\ar{لا} \ar{يتحمل} \ar{أبي} \ar{الحر}.
\end{quote}

\begin{quote}
\emph{lā yataḥammalu abī al-ḥarr.}
\end{quote}

\begin{quote}
My father cannot stand the heat.
\end{quote}

Present: \ar{يتحمل}, \emph{yataḥammal}.

\begin{cousinweb}
\textbf{\ar{تحمل}} (\emph{taḥammala}) is built on the native Arabic root \emph{ḥ-m-l}, \textquotedblleft{}to carry\textquotedblright{} (Lane). Ḥamala, to carry, which you know, is from the same root; Form V is to carry a weight oneself.
\end{cousinweb}

\begin{grammarlens}[title={What you've built}]
One more word for describing what you lived and felt.
\end{grammarlens}

\subsection*{Your turn}
\begin{itemize}
\raggedright
  \item \emph{Say it:} \emph{taḥammala}
  \item \emph{Say it:} \emph{taḥammala}, once more
  \item \emph{Say it:} \emph{lā yataḥammalu abī al-ḥarr.} — My father cannot stand the heat.
  \item \emph{Recall:} say the Arabic for annoyed, then the Arabic for displeased, then say \emph{taḥammala} again
\end{itemize}

\begin{tcolorbox}[breakable,colback=teal!4,colframe=teal!35!black,title={Before you move on}]
What does \ar{تحمل} mean? (\textquotedblleft{}To put up with\textquotedblright{}.) Say it once more.
\end{tcolorbox}
\section[\texorpdfstring{\ar{ضجر} (ḍajira)}{ضجر (ḍajira)}]{\ar{ضجر} (ḍajira) — to be fed up}

\label{lesson:AR-C495-dajira}



\begin{quote}
Before the new one: say the Arabic for displeased, then the Arabic for to put up with.
\end{quote}

\subsection*{You'll want to know: \ar{ضجر}}
\textbf{\ar{ضجر}} — \emph{ḍajira} — \textquotedblleft{}to be fed up\textquotedblright{}.

\begin{quote}
\ar{ضجر} \ar{المسافرون} \ar{من} \ar{التأخير} \ar{الطويل}.
\end{quote}

\begin{quote}
\emph{ḍajira al-musāfirūn min at-taʾkhīr aṭ-ṭawīl.}
\end{quote}

\begin{quote}
The travellers were fed up with the long delay.
\end{quote}

Present: \ar{يضجر}, \emph{yaḍjar}.

It takes \ar{من}, as \ar{مل} does.

\begin{cousinweb}
\textbf{\ar{ضجر}} (\emph{ḍajira}) is built on the native Arabic root \emph{ḍ-j-r}, \textquotedblleft{}to be vexed, disquieted\textquotedblright{} (Lane). Lane gives yaḍjaru as its present.
\end{cousinweb}

\begin{grammarlens}[title={What you've built}]
One more word for describing what you lived and felt.
\end{grammarlens}

\subsection*{Your turn}
\begin{itemize}
\raggedright
  \item \emph{Say it:} \emph{ḍajira}
  \item \emph{Say it:} \emph{ḍajira}, once more
  \item \emph{Say it:} \emph{ḍajira al-musāfirūn min at-taʾkhīr aṭ-ṭawīl.} — The travellers were fed up with the long delay.
  \item \emph{Recall:} say the Arabic for displeased, then the Arabic for to put up with, then say \emph{ḍajira} again
\end{itemize}

\begin{tcolorbox}[breakable,colback=teal!4,colframe=teal!35!black,title={Before you move on}]
What does \ar{ضجر} mean? (\textquotedblleft{}To be fed up\textquotedblright{}.) Say it once more.
\end{tcolorbox}
\section[(dialogue)]{Feelings and moods — a review}

\label{lesson:AR-R495-review-488-495}



\begin{quote}
Say the Arabic for happiness and boredom.
\end{quote}

\subsection*{Your turn}
Say these aloud:

\begin{itemize}
\raggedright
  \item happiness, boredom, to get angry, to grieve, to sigh
  \item panic, terror, calm, to shiver, to frighten
  \item astonishment, astonished, to be amazed, to be surprised, admiration
  \item proud, to take pride, gratitude, grateful, to impress
  \item jealousy, jealous, envious, to envy, to be jealous
  \item to regret, regretful, pity, to reproach, despair
  \item gloom, nostalgia, frustration, frustrated, to get bored
  \item annoyance, annoyed, displeased, to put up with, to be fed up
\end{itemize}

\begin{tcolorbox}[breakable,colback=teal!4,colframe=teal!35!black,title={Before you move on}]
Happiness? (\textbf{\ar{سعادة}}, \emph{saʿāda}.) Boredom? (\textbf{\ar{ملل}}, \emph{malal}.) To get angry? (\textbf{\ar{غضب}}, \emph{ghaḍiba}.) To grieve? (\textbf{\ar{حزن}}, \emph{ḥazina}.) To sigh? (\textbf{\ar{تنهد}}, \emph{tanahhada}.) Panic? (\textbf{\ar{ذعر}}, \emph{dhuʿr}.) Terror? (\textbf{\ar{رعب}}, \emph{ruʿb}.) Calm? (\textbf{\ar{هدوء}}, \emph{hudūʾ}.) To shiver? (\textbf{\ar{ارتعش}}, \emph{irtaʿasha}.) To frighten? (\textbf{\ar{أخاف}}, \emph{akhāfa}.) Astonishment? (\textbf{\ar{دهشة}}, \emph{dahsha}.) Astonished? (\textbf{\ar{مندهش}}, \emph{mundahish}.) To be amazed? (\textbf{\ar{اندهش}}, \emph{indahasha}.) To be surprised? (\textbf{\ar{تفاجأ}}, \emph{tafājaʾa}.) Admiration? (\textbf{\ar{إعجاب}}, \emph{iʿjāb}.) Proud? (\textbf{\ar{فخور}}, \emph{fakhūr}.) To take pride? (\textbf{\ar{افتخر}}, \emph{iftakhara}.) Gratitude? (\textbf{\ar{امتنان}}, \emph{imtinān}.) Grateful? (\textbf{\ar{ممتن}}, \emph{mumtann}.) To impress? (\textbf{\ar{أعجب}}, \emph{aʿjaba}.) Jealousy? (\textbf{\ar{غيرة}}, \emph{ghīra}.) Jealous? (\textbf{\ar{غيور}}, \emph{ghayūr}.) Envious? (\textbf{\ar{حسود}}, \emph{ḥasūd}.) To envy? (\textbf{\ar{حسد}}, \emph{ḥasada}.) To be jealous? (\textbf{\ar{غار}}, \emph{ghāra}.) To regret? (\textbf{\ar{ندم}}, \emph{nadima}.) Regretful? (\textbf{\ar{نادم}}, \emph{nādim}.) Pity? (\textbf{\ar{شفقة}}, \emph{shafaqa}.) To reproach? (\textbf{\ar{عاتب}}, \emph{ʿātaba}.) Despair? (\textbf{\ar{يأس}}, \emph{yaʾs}.) Gloom? (\textbf{\ar{كآبة}}, \emph{kaʾāba}.) Nostalgia? (\textbf{\ar{حنين}}, \emph{ḥanīn}.) Frustration? (\textbf{\ar{إحباط}}, \emph{iḥbāṭ}.) Frustrated? (\textbf{\ar{محبط}}, \emph{muḥbaṭ}.) To get bored? (\textbf{\ar{مل}}, \emph{malla}.) Annoyance? (\textbf{\ar{انزعاج}}, \emph{inziʿāj}.) Annoyed? (\textbf{\ar{منزعج}}, \emph{munzaʿij}.) Displeased? (\textbf{\ar{مستاء}}, \emph{mustāʾ}.) To put up with? (\textbf{\ar{تحمل}}, \emph{taḥammala}.) To be fed up? (\textbf{\ar{ضجر}}, \emph{ḍajira}.)
\end{tcolorbox}
